\section{Zadanie 2}
\autor{Hubert Jastrzębski}
Niech X, Y, Z będą niezależnymi zmiennymi losowymi ze standardowym rozkładem normalnym. Która z wartości P((X + Y + Z)/3 < 10) i P(X < 10) jest większa?
Odpowiedź uzasadnij.

\subsection*{Rozwiązanie}
1. $P(X < 10) = \Phi(10)$

$X + Y + Z \sim N(0, 3)$ (lemat 1).

2. $P((X + Y + Z)/3 < 10) = P((X + Y + Z) < 30) = \Phi(\frac{30}{\sqrt{3}}) = \Phi(10 \sqrt{3})$ (lemat 2).

$\Phi(10 \sqrt{3}) > \Phi(10)$, więc $P((X + Y + Z)/3 < 10) > P(X < 10)$.

Nie pamiętam czy to było na wykładzie (pewnie tak, ale \textbf{ktoś} go nie spisał, a mi się nie chce robić notatek), ale mogę przy okazji udowodnić dwie rzeczy z których korzystam,
bo samo zadanie jest dość proste:

\textbf{Lemat 1:} Jeżeli $X \sim N(\mu_1, \sigma_1^2)$ i $Y \sim N(\mu_2, \sigma_2^2)$ są niezależne, to $X + Y \sim N(\mu_1 + \mu_2, \sigma_1^2 + \sigma_2^2)$

\textbf{Dowód:}

Wiemy, że gęstości $X$ i $Y$ to:
$$f_X(x) = \frac{1}{\sqrt{2 \pi} \sigma_1} e^{-\frac{(x - \mu_1)^2}{2 \sigma_1^2}}$$
$$f_Y(y) = \frac{1}{\sqrt{2 \pi} \sigma_2} e^{-\frac{(y - \mu_2)^2}{2 \sigma_2^2}}$$

Gęstość sumy dwóch niezależnych zmiennych losowych to splot ich gęstości:
$$f_{X+Y}(z) = \int_{-\infty}^{\infty} f_X(x) f_Y(z - x) dx$$

\[f_{X + Y}(z) = \int_{-\infty}^{\infty} \frac{1}{2 \pi \sigma_1 \sigma_2} e^{-\left(\frac{(x - \mu_1)^2}{2 \sigma_1^2} + \frac{(z - x - \mu_2)^2}{2 \sigma_2^2}\right)} dx\]

Możemy wykazać, że: $\frac{(x - A)^2}{a} + \frac{(x - B)^2}{b} = \frac{a+b}{ab}\left(x - \frac{Ab + Ba}{a+b}\right)^2 + \frac{(A-B)^2}{a+b}$

Podstawiając odpowiednio dostajemy: $\frac{(x - \mu_1)^2}{2 \sigma_1^2} + \frac{(z - x - \mu_2)^2}{2 \sigma_2^2} = \frac{\sigma_1^2 + \sigma_2^2}{\sigma_1^2 \sigma_2^2} \left( x - \tilde{\mu} \right)^2 + \frac{(z - (\mu_1 + \mu_2))^2}{\sigma_1^2 + \sigma_2^2}$

\[f_{X+Y}(z) = \underbrace{\frac{1}{2\pi \sigma_1 \sigma_2} e^{-\frac{(z - (\mu_1+\mu_2))^2}{2(\sigma_1^2 + \sigma_2^2)}}}_{\text{stała względem } x} \cdot \int_{-\infty}^{\infty} e^{-\frac{1}{2} \frac{\sigma_1^2+\sigma_2^2}{\sigma_1^2\sigma_2^2}(x - \tilde{\mu})^2} dx\]

Mamy jakiś taki wzór (oczywiście każdy go zna z analizy): $\int_{-\infty}^{\infty} e^{-a(x-b)^2} dx = \sqrt{\frac{\pi}{a}}$

Korzystając z niego dostajemy: 
\[\int_{-\infty}^{\infty} e^{-\frac{1}{2} \frac{\sigma_1^2+\sigma_2^2}{\sigma_1^2\sigma_2^2}(x - \tilde{\mu})^2} dx = \sqrt{\frac{2 \pi \sigma_1^2 \sigma_2^2}{\sigma_1^2 + \sigma_2^2}}\]

Ostatecznie więc:
\[f_{X+Y}(z) = \frac{1}{\sqrt{2 \pi} \sqrt{\sigma_1^2 + \sigma_2^2}} e^{-\frac{(z - (\mu_1 + \mu_2))^2}{2(\sigma_1^2 + \sigma_2^2)}}\]

\textbf{Lemat 2:} Jeżeli $X \sim N(\mu, \sigma^2)$, to $P(X < x) = \Phi\left(\frac{x - \mu}{\sigma}\right)$

\textbf{Dowód:}

Wiemy, że:
$$\Phi(x) = \int_{-\infty}^{x} \underbrace{\frac{1}{\sqrt{2\pi}} e^{-\frac{t^2}{2}}}_{\text{gęstość } N(0, 1)} \, dt \implies \Phi\left(\frac{x - \mu}{\sigma}\right) = \int_{-\infty}^{\frac{x - \mu}{\sigma}} \frac{1}{\sqrt{2\pi}} e^{-\frac{t^2}{2}} \, dt$$
$$P(X < x) = \int_{-\infty}^{x} \underbrace{\frac{1}{\sigma\sqrt{2\pi}} e^{-\frac{(t - \mu)^2}{2\sigma^2}}}_{\text{gęstość } N(\mu, \sigma^2)} \, dt$$

Podstawmy w drugim całkowaniu $u = \frac{t - \mu}{\sigma} \implies t = \sigma u + \mu \implies dt = \sigma du$.

Granice całkowania: kiedy $t \to -\infty$, to $u \to -\infty$; kiedy $t = x$, to $u = \frac{x - \mu}{\sigma}$.
$$P(X < x) = \int_{-\infty}^{\frac{x - \mu}{\sigma}} \frac{1}{\sigma \sqrt{2\pi}} e^{-\frac{1}{2} \left(\frac{t-\mu}{\sigma}\right)^2} \sigma du = \int_{-\infty}^{\frac{x - \mu}{\sigma}} \frac{1}{\sqrt{2\pi}} e^{-\frac{u^2}{2}} \, du = \Phi\left(\frac{x - \mu}{\sigma}\right)$$


\textbf{Dodatek na koniec:}

Jeśli wiemy, że $(X + Y + Z) / 3 \sim N(0, 1/3)$, to można łatwo zrozumieć czemu $$P((X + Y + Z)/3 < 10) > P(X < 10)$$

Rozkład $N(0, 1/3)$ jest bardziej skupiony wokół średniej (mniejsza wariancja) niż rozkład $N(0, 1)$.
Zatem prawdopodobieństwo, że zmienna losowa z rozkładu $N(0, 1/3)$ przyjmie wartość mniejszą niż 10
jest większa niż prawdopodobieństwo, że zmienna losowa z rozkładu $N(0, 1)$ przyjmie wartość mniejszą niż 10.

\subsection*{Odpowiedź}
\[
\boxed{P((X + Y + Z)/3 < 10) > P(X < 10)}
\]
